% ======================================================================
% SECTION 12: TRUST, PROVENANCE, AND AUDIT INFRASTRUCTURE
% File: sections/trust_layer.tex
% ======================================================================

The trust layer provides the foundational infrastructure that enables autonomous sponsor
operation within a regulated clinical trial environment. Without verifiable trust, no
regulatory authority would accept decisions made by software agents as equivalent to
decisions made by qualified human professionals. The trust layer addresses this requirement
through four components: identity and authorization management for all agents and human
actors, immutable audit infrastructure that satisfies regulatory inspection requirements,
provenance chain architecture that traces every outcome back to its inputs and decision
logic, and a policy enforcement engine that encodes regulatory requirements as
machine-readable rules evaluated before every agent action.

\subsection{Identity and Authorization Framework}

The identity and authorization framework assigns each of the twelve functional agents a
unique cryptographic identity based on public key infrastructure (PKI). Each agent holds a
certificate issued by the trust layer's certificate authority, which establishes its identity,
authorized operations, and decision scope. Inter-agent communication uses mutual TLS
authentication: both the requesting and responding agents verify each other's certificates
before exchanging data or executing requests.

Role-based access control (RBAC) governs what each agent can do within the system. The RBAC
model defines roles corresponding to each agent type and assigns permissions that limit
each agent to operations within its functional scope. The safety agent can create and
submit safety reports but cannot modify protocol design. The clinical accountability agent
can define eligibility criteria and dose schedules but cannot authorize robotic procedures.
The robot execution gateway can issue procedure commands but cannot modify safety reporting
thresholds. This separation of privileges prevents any single agent from accumulating
sufficient authority to bypass safety controls.

Human sponsor-of-record identity management provides the authentication and authorization
infrastructure for human oversight. The sponsor-of-record holds credentials that grant
unrestricted read access to all agent activities, override authority for any agent decision,
and exclusive authority for mandatory human gate approvals. Electronic signatures by the
sponsor-of-record comply with 21 CFR Part 11 requirements \cite{cfr-part-11}: each signature
includes the signer's identity, the date and time, and the meaning of the signature.

Site-level authorization manages the tokens and sessions that govern sponsor-to-site
interactions. Each trial site receives authorization tokens that define the scope of
data exchange, robotic procedure types authorized at that site, and the personnel
authorized to interact with the sponsor system. Token lifecycle management includes
issuance upon site activation, periodic renewal, revocation upon site deactivation, and
emergency revocation when security events occur.

\subsection{Immutable Audit Infrastructure}

The immutable audit infrastructure records every agent decision, data access, state change,
and communication in a tamper-evident format that satisfies regulatory inspection
requirements. The audit system implements three layers of integrity protection.

The first layer records individual audit entries with system-generated timestamps, agent
identification, action classification, input data references, output data references,
decision rationale, and confidence score. Each entry is stored in an append-only database
that prevents modification or deletion of existing records. The entry format satisfies
21 CFR Part 11 \cite{cfr-part-11} requirements for audit trails: who did what, when, and
why, with the previous and new values for any data modification.

The second layer anchors audit entries to a blockchain or distributed ledger that provides
cryptographic proof of temporal ordering and tamper evidence. Periodic anchoring (hourly
by default, with immediate anchoring for safety-critical decisions) creates checkpoints
that an external auditor can independently verify. This layer provides evidence that
audit records have not been retroactively modified, which is a critical requirement for
regulatory inspection credibility.

The third layer implements the ICH E6(R3) essential document management requirement.
Essential documents (as defined in ICH E6(R3) Appendix) are maintained in an electronic
Trial Master File (eTMF) with lifecycle tracking from creation through archival. The
system tracks document completeness, version currency, and filing compliance, generating
alerts when essential documents are missing or outdated.

The EU Clinical Trials Regulation 25-year archiving requirement \cite{eu-ctr} is
addressed through a dedicated archive subsystem. Sponsor-appointed archive access controls
ensure that archived records remain accessible throughout the 25-year retention period.
The archive implements format migration strategies to address technology obsolescence:
records are stored in open, standardized formats (PDF/A for documents, XML for structured
data, CSV for datasets) and periodically validated for readability.

Table~\ref{tab:trust-architecture} presents the trust layer architecture.

\begin{longtable}{L{2.0cm} L{2.5cm} L{2.0cm} L{2.5cm} L{2.6cm}}
\caption{Trust Layer Architecture}
\label{tab:trust-architecture} \\
\toprule
\textbf{Layer} & \textbf{Component} & \textbf{Protocol} & \textbf{Regulatory Basis} & \textbf{Verification Method} \\
\midrule
\endfirsthead
\toprule
\textbf{Layer} & \textbf{Component} & \textbf{Protocol} & \textbf{Regulatory Basis} & \textbf{Verification Method} \\
\midrule
\endhead
\midrule
\multicolumn{5}{r}{\textit{Continued on next page}} \\
\endfoot
\bottomrule
\endlastfoot
Identity & Agent PKI certificates & X.509v3, mutual TLS & 21 CFR Part 11, ICH E6(R3) & Certificate chain validation, revocation checking \\
Identity & Human credential management & Multi-factor authentication & 21 CFR Part 11 & Identity proofing, credential lifecycle audit \\
Authorization & Role-based access control & XACML policy evaluation & ICH E6(R3) & Permission audit, least-privilege verification \\
Audit & Decision recording & Append-only database & 21 CFR Part 11 & Hash verification, completeness checks \\
Audit & Blockchain anchoring & Merkle tree proofs & Tamper evidence & Independent anchor verification \\
Audit & eTMF management & TMF Reference Model & ICH E6(R3) & Document completeness scoring \\
Provenance & Data lineage tracking & W3C PROV-DM & 21 CFR Part 11, ICH E6(R3) & End-to-end traceability verification \\
Provenance & Decision traceability & Custom provenance graph & ICH E6(R3) & Input-to-output reconstruction \\
Policy & Regulatory rule engine & Rego/OPA evaluation & Adapted 21 CFR 312/50, ICH E6(R3) & Policy coverage analysis, rule testing \\
Policy & Safety gate enforcement & State machine validation & Adapted 21 CFR 312 & Gate bypass detection, timeout monitoring \\
\end{longtable}

\subsection{Provenance Chain Architecture}

The provenance chain architecture implements the W3C PROV Data Model \cite{w3c-prov} for
tracking the lineage of every data element, decision, and physical action within the
autonomous sponsor system. Four provenance domains operate in parallel, each tracking a
different category of system activity.

Data provenance traces every data element from its original source through all
transformations to its final submission form. When a laboratory value is captured from a
site EHR system, the provenance chain records the source system, extraction timestamp,
FHIR resource identifier, transformation steps (unit conversion, normal range
classification), EDC storage location, SDTM mapping, ADaM derivation, and submission
dataset inclusion. This end-to-end traceability enables auditors to verify any submitted
data point back to its original source.

Decision provenance traces every agent decision from its triggering event through the
inputs considered, the algorithm applied, the confidence assessment, and the resulting
action. When the safety agent classifies an adverse event, the provenance chain records
the source report, the MedDRA coding decision, the causality assessment inputs and
method, the seriousness determination, the expectedness classification, and the resulting
regulatory reporting action. This provenance chain enables retrospective review of any
safety decision.

Procedure provenance reconstructs the complete timeline of every robotic clinical procedure.
This domain integrates robotic telemetry data (sensor readings at millisecond intervals),
patient monitoring data (vital signs at second intervals), clinical observations (human
documentation at event level), and system events (safety gate approvals, state transitions,
alarms) into a unified procedure record. The procedure provenance chain enables
reconstruction of exactly what happened during any robotic procedure, which is essential
for safety event investigation and regulatory review.

Model provenance tracks the lifecycle of machine learning models used within the autonomous
sponsor system. For each model, the provenance chain records the training data (with
privacy-preserving references to federated data sources), training parameters, validation
methodology, performance metrics, deployment version, and inference log. This provenance
domain addresses the emerging regulatory expectation that AI/ML systems used in clinical
trials maintain complete model lifecycle documentation.

\subsection{Policy Enforcement and Safety Gating}

The policy enforcement engine encodes the adapted regulatory frameworks (21 CFR Part 312
\cite{cfr312-adapt}, 21 CFR Part 50 \cite{cfr50-adapt}, ICH E6(R3) \cite{ich-adapt}) as
machine-readable policy rules that are evaluated before every agent action. The engine uses
Open Policy Agent (OPA) with Rego policy language to define regulatory constraints as
declarative rules that can be independently tested, versioned, and audited.

Table~\ref{tab:safety-gate-classification} presents the safety gate classification.

\begin{longtable}{L{1.8cm} L{2.5cm} L{3.0cm} L{1.5cm} L{2.8cm}}
\caption{Safety Gate Classification}
\label{tab:safety-gate-classification} \\
\toprule
\textbf{Gate Type} & \textbf{Trigger Condition} & \textbf{Required Approvals} & \textbf{Timeout} & \textbf{Default Action} \\
\midrule
\endfirsthead
\toprule
\textbf{Gate Type} & \textbf{Trigger Condition} & \textbf{Required Approvals} & \textbf{Timeout} & \textbf{Default Action} \\
\midrule
\endhead
\midrule
\multicolumn{5}{r}{\textit{Continued on next page}} \\
\endfoot
\bottomrule
\endlastfoot
Informational & Non-critical agent decision within normal parameters & None (automated with logging) & None & Proceed with audit entry \\
Confirmation & Agent decision at scope boundary or elevated confidence & Agent self-check plus peer agent validation & 5 minutes & Hold pending validation \\
Clinical & Decision affecting patient treatment or safety classification & Supervising clinician electronic approval & 30 minutes & Escalate to medical monitor \\
Regulatory & Decision affecting regulatory submissions or authority communications & Sponsor-of-record electronic signature & 4 hours & Hold pending approval \\
Safety-critical & Robotic Tier 3 procedure, SAE escalation, or emergency response & Mandatory human approval with no bypass & No timeout (wait indefinitely) & No autonomous action permitted \\
\end{longtable}

The policy engine operates on a deny-by-default principle: every agent action must be
explicitly permitted by at least one policy rule, and no action proceeds if any applicable
deny rule matches. This conservative approach ensures that regulatory requirements are
enforced even when the system encounters situations not anticipated during policy
development. When an action is denied by policy, the agent receives a structured denial
response including the specific rule that triggered the denial, enabling targeted remediation.

Continuous validation extends the initial system validation into ongoing operation. The
trust layer monitors system behavior against expected patterns and flags anomalies that
may indicate system degradation, configuration drift, or security events. Validation
evidence is generated continuously and compiled into periodic validation status reports
that contribute to the quality agent's audit readiness assessment.
